\documentclass[12pt]{article}

% Idioma y codificación
\usepackage[spanish, es-tabla]{babel}       %es-tabla para que se titule "Tabla"
\usepackage[utf8]{inputenc}

% Márgenes
\usepackage[a4paper,top=3cm,bottom=2.5cm,left=3cm,right=3cm]{geometry}

% Comentarios de bloque
\usepackage{verbatim}

% Paquetes de links
\usepackage[hidelinks]{hyperref}    % Permite enlaces
\usepackage{url}                    % redirecciona a la web

% Más opciones para enumeraciones
\usepackage{enumitem}

% Personalizar la portada
\usepackage{titling}


% Paquetes de tablas
\usepackage{multirow}


%------------------------------------------------------------------------

%Paquetes de figuras
\usepackage{caption}
\usepackage{subcaption} % Figuras al lado de otras
\usepackage{float}      % Poner figuras en el sitio indicado H.


% Paquetes de imágenes
\usepackage{graphicx}       % Paquete para añadir imágenes
\usepackage{transparent}    % Para manejar la opacidad de las figuras

% Paquete para usar colores
\usepackage[dvipsnames]{xcolor}
\usepackage{pagecolor}      % Para cambiar el color de la página

% Habilita tamaños de fuente mayores
\usepackage{fix-cm}


%------------------------------------------------------------------------

% Paquetes de matemáticas
\usepackage{mathtools, amsfonts, amssymb, mathrsfs}
\usepackage[makeroom]{cancel}     % Simplificar tachando
\usepackage{polynom}    % Divisiones y Ruffini
\usepackage{units} % Para poner fracciones diagonales con \nicefrac

\usepackage{pgfplots}   %Representar funciones
\pgfplotsset{compat=1.18}  % Versión 1.18

\usepackage{tikz-cd}    % Para usar diagramas de composiciones
\usetikzlibrary{calc}   % Para usar cálculo de coordenadas en tikz

%Definición de teoremas, etc.
\usepackage{amsthm}
%\swapnumbers   % Intercambia la posición del texto y de la numeración

\theoremstyle{plain}

\makeatletter
\@ifclassloaded{article}{
  \newtheorem{teo}{Teorema}[section]
}{
  \newtheorem{teo}{Teorema}[chapter]  % Se resetea en cada chapter
}
\makeatother

\newtheorem{coro}{Corolario}[teo]           % Se resetea en cada teorema
\newtheorem{prop}[teo]{Proposición}         % Usa el mismo contador que teorema
\newtheorem{lema}[teo]{Lema}                % Usa el mismo contador que teorema

\theoremstyle{remark}
\newtheorem*{observacion}{Observación}

\theoremstyle{definition}

\makeatletter
\@ifclassloaded{article}{
  \newtheorem{definicion}{Definición} [section]     % Se resetea en cada chapter
}{
  \newtheorem{definicion}{Definición} [chapter]     % Se resetea en cada chapter
}
\makeatother

\newtheorem*{notacion}{Notación}
\newtheorem*{ejemplo}{Ejemplo}
\newtheorem*{ejercicio*}{Ejercicio}             % No numerado
\newtheorem{ejercicio}{Ejercicio} [section]     % Se resetea en cada section


% Modificar el formato de la numeración del teorema "ejercicio"
\renewcommand{\theejercicio}{%
  \ifnum\value{section}=0 % Si no se ha iniciado ninguna sección
    \arabic{ejercicio}% Solo mostrar el número de ejercicio
  \else
    \thesection.\arabic{ejercicio}% Mostrar número de sección y número de ejercicio
  \fi
}


% \renewcommand\qedsymbol{$\blacksquare$}         % Cambiar símbolo QED
%------------------------------------------------------------------------

% Paquetes para encabezados
\usepackage{fancyhdr}
\pagestyle{fancy}
\fancyhf{}

\newcommand{\helv}{ % Modificación tamaño de letra
\fontfamily{}\fontsize{12}{12}\selectfont}
\setlength{\headheight}{15pt} % Amplía el tamaño del índice


%\usepackage{lastpage}   % Referenciar última pag   \pageref{LastPage}
\fancyfoot[C]{\thepage}

%------------------------------------------------------------------------

% Conseguir que no ponga "Capítulo 1". Sino solo "1."
\makeatletter
\@ifclassloaded{book}{
  \renewcommand{\chaptermark}[1]{\markboth{\thechapter.\ #1}{}} % En el encabezado
    
  \renewcommand{\@makechapterhead}[1]{%
  \vspace*{50\p@}%
  {\parindent \z@ \raggedright \normalfont
    \ifnum \c@secnumdepth >\m@ne
      \huge\bfseries \thechapter.\hspace{1em}\ignorespaces
    \fi
    \interlinepenalty\@M
    \Huge \bfseries #1\par\nobreak
    \vskip 40\p@
  }}
}
\makeatother

%------------------------------------------------------------------------
% Paquetes de cógido
\usepackage{minted}
\renewcommand\listingscaption{Código fuente}

\usepackage{fancyvrb}
% Personaliza el tamaño de los números de línea
\renewcommand{\theFancyVerbLine}{\small\arabic{FancyVerbLine}}

% Estilo para C++
\newminted{cpp}{
    frame=lines,
    framesep=2mm,
    baselinestretch=1.2,
    linenos,
    escapeinside=||
}



\usepackage{listings} % Para incluir código desde un archivo

\renewcommand\lstlistingname{Código Fuente}
\renewcommand\lstlistlistingname{Índice de Códigos Fuente}

% Definir colores
\definecolor{vscodepurple}{rgb}{0.5,0,0.5}
\definecolor{vscodeblue}{rgb}{0,0,0.8}
\definecolor{vscodegreen}{rgb}{0,0.5,0}
\definecolor{vscodegray}{rgb}{0.5,0.5,0.5}
\definecolor{vscodebackground}{rgb}{0.97,0.97,0.97}
\definecolor{vscodelightgray}{rgb}{0.9,0.9,0.9}

% Configuración para el estilo de C similar a VSCode
\lstdefinestyle{vscode_C}{
  backgroundcolor=\color{vscodebackground},
  commentstyle=\color{vscodegreen},
  keywordstyle=\color{vscodeblue},
  numberstyle=\tiny\color{vscodegray},
  stringstyle=\color{vscodepurple},
  basicstyle=\scriptsize\ttfamily,
  breakatwhitespace=false,
  breaklines=true,
  captionpos=b,
  keepspaces=true,
  numbers=left,
  numbersep=5pt,
  showspaces=false,
  showstringspaces=false,
  showtabs=false,
  tabsize=2,
  frame=tb,
  framerule=0pt,
  aboveskip=10pt,
  belowskip=10pt,
  xleftmargin=10pt,
  xrightmargin=10pt,
  framexleftmargin=10pt,
  framexrightmargin=10pt,
  framesep=0pt,
  rulecolor=\color{vscodelightgray},
  backgroundcolor=\color{vscodebackground},
}

%------------------------------------------------------------------------

% Comandos definidos
\newcommand{\bb}[1]{\mathbb{#1}}
\newcommand{\cc}[1]{\mathcal{#1}}

% I prefer the slanted \leq
\let\oldleq\leq % save them in case they're every wanted
\let\oldgeq\geq
\renewcommand{\leq}{\leqslant}
\renewcommand{\geq}{\geqslant}

% Si y solo si
\newcommand{\sii}{\iff}

% Letras griegas
\newcommand{\eps}{\epsilon}
\newcommand{\veps}{\varepsilon}
\newcommand{\lm}{\lambda}

\newcommand{\ol}{\overline}
\newcommand{\ul}{\underline}
\newcommand{\wt}{\widetilde}
\newcommand{\wh}{\widehat}

\let\oldvec\vec
\renewcommand{\vec}{\overrightarrow}

% Derivadas parciales
\newcommand{\del}[2]{\frac{\partial #1}{\partial #2}}
\newcommand{\Del}[3]{\frac{\partial^{#1} #2}{\partial^{#1} #3}}
\newcommand{\deld}[2]{\dfrac{\partial #1}{\partial #2}}
\newcommand{\Deld}[3]{\dfrac{\partial^{#1} #2}{\partial^{#1} #3}}


\newcommand{\AstIg}{\stackrel{(\ast)}{=}}
\newcommand{\Hop}{\stackrel{L'H\hat{o}pital}{=}}

\newcommand{\red}[1]{{\color{red}#1}} % Para integrales, destacar los cambios.

% Método de integración
\newcommand{\MetInt}[2]{
    \left[\begin{array}{c}
        #1 \\ #2
    \end{array}\right]
}

% Declarar aplicaciones
% 1. Nombre aplicación
% 2. Dominio
% 3. Codominio
% 4. Variable
% 5. Imagen de la variable
\newcommand{\Func}[5]{
    \begin{equation*}
        \begin{array}{rrll}
            #1:& #2 & \longrightarrow & #3\\
               & #4 & \longmapsto & #5
        \end{array}
    \end{equation*}
}

%------------------------------------------------------------------------
\usepgfplotslibrary{fillbetween}

\begin{document}

    % 1. Foto de fondo
    % 2. Título
    % 3. Encabezado Izquierdo
    % 4. Color de fondo
    % 5. Coord x del titulo
    % 6. Coord y del titulo
    % 7. Fecha

    
    % 1. Foto de fondo
% 2. Título
% 3. Encabezado Izquierdo
% 4. Color de fondo
% 5. Coord x del titulo
% 6. Coord y del titulo
% 7. Fecha

\newcommand{\portada}[7]{

    \portadaBase{#1}{#2}{#3}{#4}{#5}{#6}{#7}
    \portadaBook{#1}{#2}{#3}{#4}{#5}{#6}{#7}
}

\newcommand{\portadaExamen}[7]{

    \portadaBase{#1}{#2}{#3}{#4}{#5}{#6}{#7}
    \portadaArticle{#1}{#2}{#3}{#4}{#5}{#6}{#7}
}




\newcommand{\portadaBase}[7]{

    % Tiene la portada principal y la licencia Creative Commons
    
    % 1. Foto de fondo
    % 2. Título
    % 3. Encabezado Izquierdo
    % 4. Color de fondo
    % 5. Coord x del titulo
    % 6. Coord y del titulo
    % 7. Fecha
    
    
    \thispagestyle{empty}               % Sin encabezado ni pie de página
    \newgeometry{margin=0cm}        % Márgenes nulos para la primera página
    
    
    % Encabezado
    \fancyhead[L]{\helv #3}
    \fancyhead[R]{\helv \nouppercase{\leftmark}}
    
    
    \pagecolor{#4}        % Color de fondo para la portada
    
    \begin{figure}[p]
        \centering
        \transparent{0.3}           % Opacidad del 30% para la imagen
        
        \includegraphics[width=\paperwidth, keepaspectratio]{assets/#1}
    
        \begin{tikzpicture}[remember picture, overlay]
            \node[anchor=north west, text=white, opacity=1, font=\fontsize{60}{90}\selectfont\bfseries\sffamily, align=left] at (#5, #6) {#2};
            
            \node[anchor=south east, text=white, opacity=1, font=\fontsize{12}{18}\selectfont\sffamily, align=right] at (9.7, 3) {\textbf{\href{https://losdeldgiim.github.io/}{Los Del DGIIM}}};
            
            \node[anchor=south east, text=white, opacity=1, font=\fontsize{12}{15}\selectfont\sffamily, align=right] at (9.7, 1.8) {Doble Grado en Ingeniería Informática y Matemáticas\\Universidad de Granada};
        \end{tikzpicture}
    \end{figure}
    
    
    \restoregeometry        % Restaurar márgenes normales para las páginas subsiguientes
    \pagecolor{white}       % Restaurar el color de página
    
    
    \newpage
    \thispagestyle{empty}               % Sin encabezado ni pie de página
    \begin{tikzpicture}[remember picture, overlay]
        \node[anchor=south west, inner sep=3cm] at (current page.south west) {
            \begin{minipage}{0.5\paperwidth}
                \href{https://creativecommons.org/licenses/by-nc-nd/4.0/}{
                    \includegraphics[height=2cm]{assets/Licencia.png}
                }\vspace{1cm}\\
                Esta obra está bajo una
                \href{https://creativecommons.org/licenses/by-nc-nd/4.0/}{
                    Licencia Creative Commons Atribución-NoComercial-SinDerivadas 4.0 Internacional (CC BY-NC-ND 4.0).
                }\\
    
                Eres libre de compartir y redistribuir el contenido de esta obra en cualquier medio o formato, siempre y cuando des el crédito adecuado a los autores originales y no persigas fines comerciales. 
            \end{minipage}
        };
    \end{tikzpicture}
    
    
    
    % 1. Foto de fondo
    % 2. Título
    % 3. Encabezado Izquierdo
    % 4. Color de fondo
    % 5. Coord x del titulo
    % 6. Coord y del titulo
    % 7. Fecha


}


\newcommand{\portadaBook}[7]{

    % 1. Foto de fondo
    % 2. Título
    % 3. Encabezado Izquierdo
    % 4. Color de fondo
    % 5. Coord x del titulo
    % 6. Coord y del titulo
    % 7. Fecha

    % Personaliza el formato del título
    \pretitle{\begin{center}\bfseries\fontsize{42}{56}\selectfont}
    \posttitle{\par\end{center}\vspace{2em}}
    
    % Personaliza el formato del autor
    \preauthor{\begin{center}\Large}
    \postauthor{\par\end{center}\vfill}
    
    % Personaliza el formato de la fecha
    \predate{\begin{center}\huge}
    \postdate{\par\end{center}\vspace{2em}}
    
    \title{#2}
    \author{\href{https://losdeldgiim.github.io/}{Los Del DGIIM}}
    \date{Granada, #7}
    \maketitle
    
    \tableofcontents
}




\newcommand{\portadaArticle}[7]{

    % 1. Foto de fondo
    % 2. Título
    % 3. Encabezado Izquierdo
    % 4. Color de fondo
    % 5. Coord x del titulo
    % 6. Coord y del titulo
    % 7. Fecha

    % Personaliza el formato del título
    \pretitle{\begin{center}\bfseries\fontsize{42}{56}\selectfont}
    \posttitle{\par\end{center}\vspace{2em}}
    
    % Personaliza el formato del autor
    \preauthor{\begin{center}\Large}
    \postauthor{\par\end{center}\vspace{3em}}
    
    % Personaliza el formato de la fecha
    \predate{\begin{center}\huge}
    \postdate{\par\end{center}\vspace{5em}}
    
    \title{#2}
    \author{\href{https://losdeldgiim.github.io/}{Los Del DGIIM}}
    \date{Granada, #7}
    \thispagestyle{empty}               % Sin encabezado ni pie de página
    \maketitle
    \vfill
}
    \portadaExamen{ffccA4.jpg}{Probabilidad\\Examen I}{Probabilidad. Examen I}{MidnightBlue}{-8}{28}{2024-2025}{Arturo Olivares Martos}

    \begin{description}
        \item[Asignatura] Probabilidad.
        \item[Curso Académico] 2021-22.
        \item[Grado] Doble Grado en Ingeniería Informática y Matemáticas.
        \item[Grupo] Único.
        %\item[Profesor] José María Espinar García.
        \item[Descripción] Parcial de los Temas 1 y 2.1.
        \item[Fecha] 8 de noviembre de 2021.
        % \item[Duración] 60 minutos.
    
    \end{description}
    \newpage

    \begin{ejercicio}
        Dado el vector bidimensional $(X,Y)$ distribuido uniformemente en el recinto limitado
        $$R = \{(x, y) \in \mathbb{R}^2 \mid -2 \leq x \leq y \leq -x\}$$

        \begin{enumerate}
            \item ($0.5$ puntos) Obtener su función de densidad conjunta.
            
            Representamos en primer lugar dicho conjunto:
            \begin{figure}[H]
                \centering
                \begin{tikzpicture}
                    \begin{axis}
                        [
                            axis lines = center,
                            xlabel = $X$,
                            ylabel = $Y$,
                            xmin = -3, xmax = 3,
                            ymin = -3, ymax = 3,
                        ]
                        \addplot[fill=red, opacity=0.3] coordinates {(-2, -2) (0, 0) (-2, 2)} -- cycle;
                        \node at (-1.5, 0.5) {$R$};

                        % R0: (-2,3) (-2,-2) (3,-2) (3,-3) (-3,-3) (-3,3) (-2,3)
                        \addplot[fill=blue, opacity=0.2] coordinates {
                            (-2, 3) (-2, -2) (3, -2) (3, -3) (-3, -3) (-3, 3) (-2, 3)
                        } -- cycle;
                        \node at (-2.5, -2.5) {$R_0$};

                        % R1: (-2,-2) (0,0) (3,0) (3,-2)
                        \addplot[fill=green, opacity=0.2] coordinates {
                            (-2, -2) (0, 0) (3, 0) (3, -2)
                        } -- cycle;
                        \node at (1.5, -1) {$R_1$};

                        % R2: (-2,2) (0,0) (0,2)
                        \addplot[fill=yellow, opacity=0.2] coordinates {
                            (-2, 2) (0, 0) (0, 2)
                        } -- cycle;
                        \node at (-0.8, 1.3) {$R_2$};

                        % R3: (0,2) (0,0) (3,0) (3,2)
                        \addplot[fill=orange, opacity=0.2] coordinates {
                            (0, 2) (0, 0) (3, 0) (3, 2)
                        } -- cycle;
                        \node at (1.5, 1) {$R_3$};

                        % R4: (-2,2) (0,2) (0,3) (-2,3)
                        \addplot[fill=purple, opacity=0.2] coordinates {
                            (-2, 2) (0, 2) (0, 3) (-2, 3)
                        } -- cycle;
                        \node at (-1, 2.5) {$R_4$};

                        % R5: (0,2) (3,2) (3,3) (0,3)
                        \addplot[fill=gray, opacity=0.2] coordinates {
                            (0, 2) (3, 2) (3, 3) (0, 3)
                        } -- cycle;
                        \node at (1.5, 2.5) {$R_5$};
                    \end{axis}
                \end{tikzpicture}
            \end{figure}

            Tenemos que, para $x\in [-2, 0]$ y $y\in [x, -x]$, la función de densidad conjunta es:
            \begin{equation*}
                f_{(X,Y)}(x, y) = k,\qquad k\in \bb{R}^+
            \end{equation*}

            Para que sea una función de densidad, hemos de tener que:
            \begin{equation*}
                1 = \int_{\bb{R}^2} f = \int_R f = \int_R k
            \end{equation*}

            Tenemos dos opciones:
            \begin{description}
                \item[Integración Normal:]
                \begin{align*}
                    1 = \int_R k &= \int_{-2}^0 \int_{x}^{-x} k \, dy \, dx = \int_{-2}^0 k(-2x) \, dx =\\&= -2k\left[\dfrac{x^2}{2}\right]_{-2}^0 = -2k\left[0-2\right] = 4k \Longrightarrow k = \dfrac{1}{4}
                \end{align*}

                \item[Razonando según la Forma de $R$:]
                \begin{align*}
                    1 = \int_R k = k\lm(R) = k\cdot \dfrac{4\cdot 2}{2} = 4k \Longrightarrow k = \dfrac{1}{4}
                \end{align*}
            \end{description}
            En cualquier caso, para el valor de $k=\nicefrac{1}{4}$, la función de densidad conjunta integrable, no negativa e integra $1$.
            \item ($1.5$ puntos) Obtener su función de distribución conjunta.
            
            Distinguimos casos:
            \begin{itemize}
                \item \ul{Si $x < -2$ o $y < -2$} (Zona $R_0$):
                \begin{equation*}
                    F_{(X,Y)}(x, y) \int_{-\infty}^x \int_{-\infty}^y f_{(X,Y)}(u, v) \, dv \, du = 0
                \end{equation*}

                \item \ul{Si $x\in [-2, 0]$ y $y\in [x, -x]$} (Zona $R$):
                \begin{align*}
                    F_{(X,Y)}(x, y) &= \int_{-\infty}^x \int_{-\infty}^y f_{(X,Y)}(u, v) \, dv \, du = \int_{-2}^x \int_{u}^{y} \dfrac{1}{4} \, dv \, du
                    =\\&= \int_{-2}^x \dfrac{1}{4}(y-u) \, du = \dfrac{1}{4}\left[yu - \dfrac{u^2}{2}\right]_{-2}^x = \dfrac{1}{4}\left[yx - \dfrac{x^2}{2} + 2y+2\right]
                \end{align*}

                \item \ul{Si $y<0$ y $x > y$} (Zona $R_1$):
                \begin{align*}
                    F_{(X,Y)}(x, y) &= \int_{-\infty}^x \int_{-\infty}^y f_{(X,Y)}(u, v) \, dv \, du = \int_{-2}^y \int_{u}^{y} \dfrac{1}{4} \, dv \, du
                    =\\&= \int_{-2}^y \dfrac{1}{4}(y-u) \, du = \dfrac{1}{4}\left[yu - \dfrac{u^2}{2}\right]_{-2}^y = \dfrac{1}{4}\left[y^2 - \dfrac{y^2}{2} + 2y+2\right]
                    =\\&= \dfrac{1}{4}\left[\dfrac{y^2}{2} + 2y + 2\right]
                \end{align*}

                \item \ul{Si $x\in [-2, 0]$ y $y\in [-x, 2]$} (Zona $R_2$):
                \begin{align*}
                    F_{(X,Y)}(x, y) &= \int_{-\infty}^x \int_{-\infty}^y f_{(X,Y)}(u, v) \, dv \, du = \int_{-2}^{-y} \int_{u}^{y} \dfrac{1}{4} \, dv \, du+ \int_{-y}^{x} \int_{u}^{-u} \dfrac{1}{4} \, dv \, du
                    =\\&= \int_{-2}^{-y} \dfrac{1}{4}(y-u) \, du + \int_{-y}^{x} \dfrac{1}{4}(-2u) \, du = \dfrac{1}{4}\left[yu - \dfrac{u^2}{2}\right]_{-2}^{-y} + \dfrac{1}{4}\left[-u^2\right]_{-y}^{x}
                    =\\&= \dfrac{1}{4}\left[\cancel{y(-y)} - \dfrac{(-y)^2}{2} + 2y+2\right] + \dfrac{1}{4}\left[-x^2 + \cancel{y^2}\right]
                    =\\&= \dfrac{1}{4}\left[2y + 2 - \dfrac{y^2}{2} - x^2\right]
                \end{align*}

                \item \ul{Si $y\in [0, 2]$ y $x \geq 0$} (Zona $R_3$):
                \begin{align*}
                    F_{(X,Y)}(x, y) &= \int_{-\infty}^x \int_{-\infty}^y f_{(X,Y)}(u, v) \, dv \, du = \int_{-2}^{-y} \int_{u}^{y} \dfrac{1}{4} \, dv \, du+ \int_{-y}^{0} \int_{u}^{-u} \dfrac{1}{4} \, dv \, du
                    =\\&= \int_{-2}^{-y} \dfrac{1}{4}(y-u) \, du + \int_{-y}^{0} \dfrac{1}{4}(-2u) \, du = \dfrac{1}{4}\left[yu - \dfrac{u^2}{2}\right]_{-2}^{-y} + \dfrac{1}{4}\left[-u^2\right]_{-y}^{0}
                    =\\&= \dfrac{1}{4}\left[\cancel{y(-y)} - \dfrac{(-y)^2}{2} + 2y+2\right] + \dfrac{1}{4}\left[0 - \cancel{y^2}\right]
                    =\\&= \dfrac{1}{4}\left[2y + 2 - \dfrac{y^2}{2}\right]
                \end{align*}

                \item \ul{Si $x\in [-2, 0]$ y $y\geq 2$} (Zona $R_4$):
                \begin{align*}
                    F_{(X,Y)}(x, y) &= \int_{-\infty}^x \int_{-\infty}^y f_{(X,Y)}(u, v) \, dv \, du = \int_{-2}^{x} \int_{u}^{-u} \dfrac{1}{4} \, dv \, du
                    =\\&= \int_{-2}^{x} \dfrac{1}{4}(-2u) \, du = \dfrac{1}{4}\left[-u^2\right]_{-2}^{x} = \dfrac{1}{4}\left[-x^2 + 4\right]
                \end{align*}

                \item \ul{Si $x\geq 0, y\geq 2$} (Zona $R_5$):
                \begin{equation*}
                    F_{(X,Y)}(x, y) = 1
                \end{equation*}
            \end{itemize}

            Por tanto, la función de distribución conjunta es:
            \begin{equation*}
                F_{(X,Y)}(x, y) = \begin{cases}
                    0 & \text{si } x < -2 \text{ o } y < -2\\
                    \dfrac{1}{4}\left[yx - \dfrac{x^2}{2} + 2y+2\right] & \text{si } x\in [-2, 0] \text{ y } y\in [x, -x]\\
                    \dfrac{1}{4}\left[\dfrac{y^2}{2} + 2y + 2\right] & \text{si } y<0 \text{ y } x > y\\
                    \dfrac{1}{4}\left[2y + 2 - \dfrac{y^2}{2} - x^2\right] & \text{si } x\in [-2, 0] \text{ y } y\in [-x, 2]\\
                    \dfrac{1}{4}\left[2y + 2 - \dfrac{y^2}{2}\right] & \text{si } y\in [0, 2] \text{ y } x \geq 0\\
                    \dfrac{1}{4}\left[-x^2 + 4\right] & \text{si } x\in [-2, 0] \text{ y } y\geq 2\\
                    1 & \text{si } x\geq 0, y\geq 2
                \end{cases}
            \end{equation*}
            \item ($1$ punto) Obtener las dos distribuciones condicionadas posibles.
            
            En primer lugar, hemos de obtener las funciones de densidad marginales. Respecto de $X$, tenemos:
            \begin{align*}
                f_X(x) &= \int_{\bb{R}} f_{(X,Y)}(x, y) \, dy = \int_{x}^{-x} \dfrac{1}{4} \, dy = \dfrac{1}{4}(-x-x) = -\dfrac{x}{2} \qquad x\in [-2, 0]\\
            \end{align*}

            Respecto de $Y$, tenemos:
            \begin{equation*}
                f_Y(y) = \left\{\begin{array}{ll}
                    \displaystyle \int_{-2}^{y} \dfrac{1}{4} \, dx = \dfrac{1}{4}(y+2) & \text{si } y\in [-2, 0]\\\\
                    \displaystyle \int_{-2}^{-y} \dfrac{1}{4} \, dx  = \dfrac{1}{4}(-y+2) & \text{si } y\in [0, 2]
                    \end{array}\right\}
                    = \left\{\dfrac{1}{4} (-|y|+2) \qquad y\in [-2, 2]\right.
            \end{equation*}

            Por tanto, fijado $y^*\in [-2,2]$, la función de densidad condicionada de $X$ es:
            \begin{equation*}
                f_{X|Y=y^*}(x) = \dfrac{f_{(X,Y)}(x, y^*)}{f_Y(y^*)} = \dfrac{\nicefrac{1}{4}}{\nicefrac{1}{4}(-|y^*|+2)} = \dfrac{1}{-|y^*|+2} \qquad x\in [-2, -|y^*|]
            \end{equation*}

            Por otro lado, fijado $x^*\in [-2,0]$, la función de densidad condicionada de $Y$ es:
            \begin{equation*}
                f_{Y|X=x^*}(y) = \dfrac{f_{(X,Y)}(x^*, y)}{f_X(x^*)} = \dfrac{\nicefrac{1}{4}}{-\nicefrac{x^*}{2}} = -\dfrac{1}{2x^*} \qquad y\in [x^*, -x^*]
            \end{equation*}

            \item($0.5$ puntos) Obtener la probabilidad de que $X + Y + 1 \geq 0$.
            
            Veamos qué conjunto representa $X + Y + 1 \geq 0$:
            \begin{figure}[H]
                \centering
                \begin{tikzpicture}
                    \begin{axis}
                        [
                            axis lines = center,
                            xlabel = $X$,
                            ylabel = $Y$,
                            xmin = -3, xmax = 3,
                            ymin = -3, ymax = 3,
                        ]
                        \addplot[name path=A, fill=red, opacity=0.3] coordinates {(-2, -2) (0, 0) (-2, 2)} -- cycle;

                        % Recta y = -x-1
                        \addplot[name path=B, domain=-3:2.5, samples=2, dashed] {-x-1};

                        % Add a node to label the line
                        \node at (axis cs:1,-2) [anchor=west] {$y=-x-1$};

                        % R0: (-2,2) (0,0) (-0.5,-0.5) (-2, 1.5)
                        \addplot[fill=blue, opacity=0.2] coordinates {
                            (-2, 2) (0, 0) (-0.5, -0.5) (-2, 1)
                        } -- cycle;
                    \end{axis}
                \end{tikzpicture}
            \end{figure}

            Tenemos por tanto que:
            \begin{align*}
                P[X + Y + 1 \geq 0] &= \int_{-2}^{\nicefrac{-1}{2}} \int_{-x-1}^{-x} \dfrac{1}{4} \, dy \, dx + \int_{\nicefrac{-1}{2}}^{0} \int_{x}^{-x} \dfrac{1}{4} \, dy \, dx
                =\\&= \int_{-2}^{\nicefrac{-1}{2}} \dfrac{1}{4}(-x+1+x) \, dx + \int_{\nicefrac{-1}{2}}^{0} \dfrac{1}{4}(-2x) \, dx
                =\\&= \int_{-2}^{\nicefrac{-1}{2}} \dfrac{1}{4} \, dx + \int_{\nicefrac{-1}{2}}^{0} -\dfrac{1}{2}x \, dx
                =\\&= \dfrac{1}{4}\left[x\right]_{-2}^{\nicefrac{-1}{2}} - \dfrac{1}{2}\left[\dfrac{x^2}{2}\right]_{\nicefrac{-1}{2}}^{0}
                = \dfrac{1}{4}\left[\dfrac{-1}{2}+2\right] - \dfrac{1}{2}\left[0 - \dfrac{1}{8}\right] = \dfrac{7}{16}
            \end{align*}
            \item ($1.5$ puntos) Obtener la distribución marginal de $Z = X + Y$.\\
            
            Definimos la transformación:
            \Func{g}{\bb{R}^2}{\bb{R}^2}{(X,Y)}{(Z,T)=(X+Y,Y)}

            Para obtener $g^{-1}$, buscamos despejar $(X,Y)$ en función de $(Z,T)$:
            \begin{equation*}
                \left\{
                    \begin{array}{l}
                        Z=X+Y\\
                        T=Y
                    \end{array}
                \right\}
                \Longrightarrow
                \left\{
                    \begin{array}{l}
                        X=Z-T\\
                        Y=T
                    \end{array}
                \right\}
            \end{equation*}

            Por tanto, la inversa es:
            \Func{g^{-1}}{\bb{R}^2}{\bb{R}^2}{(Z,T)}{(X,Y)=(Z-T,T)}

            Tenemos que todas las componentes de $g^{-1}$ son derivables:
            \begin{align*}
                \dfrac{\partial X}{\partial Z}(Z,T)&=1, & \dfrac{\partial X}{\partial T}(Z,T)&=-1,\\
                \dfrac{\partial Y}{\partial Z}(Z,T)&=0, & \dfrac{\partial Y}{\partial T}(Z,T)&=1.
            \end{align*}

            Además, tenemos que:
            \begin{equation*}
                \det Jg^{-1}(z,t) = \begin{vmatrix}
                    1 & -1\\
                    0 & 1
                \end{vmatrix} = 1\neq 0\qquad \forall (z,t)\in \bb{R}^2
            \end{equation*}

            Por tanto, $g(Z,T)$ es un vector aleatorio. Su función de densidad conjunta es:
            \begin{align*}
                f_{(Z,T)}(z,t) &= f_{(X,Y)}(z-t,t) \cdot |Jg^{-1}(z,t)| = f_{(X,Y)}(z-t,t) = \dfrac{1}{4}\qquad (z-t,t)\in R
            \end{align*}

            Veamos cuál es este conjunto:
            \begin{align*}
                (z-t,t) \in R &\Longleftrightarrow -2\leq z-t\leq t\leq -z+t
            \end{align*}

            Gráficamente, lo vemos en la Figura~\ref{fig:conjuntoZT}.
            \begin{figure}[H]
                \centering
                \begin{tikzpicture}
                    \begin{axis}
                        [
                            axis lines = center,
                            xlabel = $Z$,
                            ylabel = $T$,
                            xmin = -4.5, xmax = 3,
                            ymin = -4.5, ymax = 3,
                            legend pos = south east,
                        ]
                        % Recta -2=z-t
                        \addplot[name path=A, domain=-5:3, samples=2, red] {2+x};
                        \addlegendentry{$-2=z-t$}

                        % Recta z-t=t
                        \addplot[name path=B, domain=-5:3, samples=2, blue] {x/2};
                        \addlegendentry{$z-t=t$}

                        % Recta x=0
                        \addplot[teal, thick] coordinates {(0,-3) (0,3)};
                        \addlegendentry{$t=-z+t$}

                        % Fill
                        \addplot[fill=orange, opacity=0.3] fill between[of=A and B,
                        soft clip={domain=-4:0}];
                        \addlegendentry{$g(R)$}

                        % Punto (-4,-2)
                        \addplot[mark=*] coordinates {(-4, -2)};
                        \node at (axis cs:-4, -2) [anchor=west] {$(-4,-2)$};
                \end{axis}
                \end{tikzpicture}
                \caption{Conjunto $g(R)=\{(z,t)\in \bb{R}^2\mid (z-t,t)\in R\}$.}
                \label{fig:conjuntoZT}
            \end{figure}

            Por tanto, la marginal de $Z$ es, para cada $z\in \left[-4, 0\right]$:
            \begin{align*}
                f_Z(z) &= \int_{\bb{R}} f_{(Z,T)}(z, t) \, dt = \int_{\nicefrac{z}{2}}^{z+2} \dfrac{1}{4} \, dt = \dfrac{1}{4}\left[t\right]_{\nicefrac{z}{2}}^{z+2} = \dfrac{1}{4}\left[z+2-\dfrac{z}{2}\right] = \dfrac{1}{4}\left[\dfrac{z}{2}+2\right]
            \end{align*}

            Por tanto, la función de densidad marginal de $Z=X+Y$ es:
            \begin{equation*}
                f_Z(z) = \begin{cases}
                    \dfrac{1}{4}\left[\dfrac{z}{2}+2\right] & \text{si } z\in \left[-4, 0\right]\\
                    0 & \text{en otro caso}
                \end{cases}
            \end{equation*}
        \end{enumerate}
    \end{ejercicio}

    \begin{ejercicio}
        Dado el vector bidimensional (X,Y) con la siguiente función masa de probabilidad conjunta:
        \begin{equation*}
            \begin{array}{|c|ccc|}
                \hline
                X\backslash Y & 0 & 1 & 2\\
                \hline
                1 & \nicefrac{1}{4} & 0 & 0\\
                2 & 0 & \nicefrac{1}{4} & 0\\
                3 & \nicefrac{1}{4} & 0 & \nicefrac{1}{4}\\
                \hline
            \end{array}
        \end{equation*}

        \begin{enumerate}
            \item ($1$ punto) Obtener la distribución condicionada de la variable aleatoria Y a todos los posibles valores de la variable aleatoria X.
            
            Tenemos que:
            \begin{equation*}
                P[X=1] = \nicefrac{1}{4},\qquad P[X=2] = \nicefrac{1}{4},\qquad P[X=3] = \nicefrac{1}{2}
            \end{equation*}

            Por tanto, tenemos que:
            \begin{align*}
                P[Y=y|X=1] &= \begin{cases}
                    \nicefrac{1}{4} & \text{si } y=0\\
                    0 & \text{si } y\neq 0
                \end{cases} \\
                P[Y=y|X=2] &= \begin{cases}
                    \nicefrac{1}{4} & \text{si } y=1 \\
                    0 & \text{si } y\neq 1
                \end{cases} \\
                P[Y=y|X=3] &= \begin{cases}
                    \nicefrac{1}{2} & \text{si } y=0,2\\
                    0 & \text{si } y=1
                \end{cases}
            \end{align*}
            \item ($1.5$ puntos) Obtener la función masa de probabilidad conjunta del vector $(X +Y,X -Y)$.
            
            Definimos la transformación:
            \Func{g}{\bb{R}^2}{\bb{R}^2}{(X,Y)}{(Z,T)=(X+Y,X-Y)}

            Para obtener $g^{-1}$, buscamos despejar $(X,Y)$ en función de $(Z,T)$:
            \begin{equation*}
                \left\{
                    \begin{array}{l}
                        Z=X+Y\\
                        T=X-Y
                    \end{array}
                \right\}
                \Longrightarrow
                \left\{
                    \begin{array}{l}
                        X=\dfrac{Z+T}{2}\\
                        Y=\dfrac{Z-T}{2}
                    \end{array}
                \right\}
            \end{equation*}

            Por tanto, la inversa es:
            \Func{g^{-1}}{\bb{R}^2}{\bb{R}^2}{(Z,T)}{(X,Y)=\left(\dfrac{Z+T}{2},\dfrac{Z-T}{2}\right)}

            Como $X$ toma valores en $E_X=\{1,2,3\}$ y $Y$ en $E_Y=\{0,1,2\}$, tenemos que:
            \begin{equation*}
                E_Z = \{1,2,3,4,5\},\qquad E_T = \{-1,0,1,2,3\}
            \end{equation*}

            Tenemos por tanto que:
            \begin{equation*}
                P[(Z,T)=(z,t)] = P[(X,Y)=g^{-1}(z,t)] \qquad \text{si } g^{-1}(z,t)\in E_X\times E_Y
            \end{equation*}

            Por tanto, la función masa de probabilidad conjunta de $(Z,T)$ es:
            \begin{equation*}
                \begin{array}{|c|ccccc|}
                    \hline
                    Z\backslash T & -1 & 0 & 1 & 2 & 3\\
                    \hline
                    1 & 0 & 0 & \nicefrac{1}{4} & 0 & 0\\
                    2 & 0 & 0 & 0 & 0 & 0 \\
                    3 & 0 & 0 & \nicefrac{1}{4} & 0 & \nicefrac{1}{4}\\
                    4 & 0 & 0 & 0 & 0 & 0\\
                    5 & 0 & 0 & \nicefrac{1}{4} & 0 & 0\\
                    \hline
                \end{array}
            \end{equation*}
            \item ($0.5$ puntos) A partir de la función masa de probabilidad marginal de $X + Y$, obtener $P[X +Y \leq 4]$.
            
            La marginal de $Z=X+Y$ es:
            \begin{align*}
                P[Z=1] &= \nicefrac{1}{4}\\
                P[Z=2] &= 0\\
                P[Z=3] &= \nicefrac{1}{2}\\
                P[Z=4] &= 0\\
                P[Z=5] &= \nicefrac{1}{4}
            \end{align*}

            Por tanto, la probabilidad de que $X+Y\leq 4$ es:
            \begin{equation*}
                P[X+Y\leq 4] = P[Z=1] + P[Z=2] + P[Z=3] + P[Z=4] = \frac{3}{4}
            \end{equation*}
        \end{enumerate}
    \end{ejercicio}

    \begin{ejercicio}
        Realizar los siguientes apartados:
        \begin{enumerate}
            \item ($1$ punto) Deducir de forma razonada la función generatriz de momentos de una distribución Gamma. ¿Está definida para todos los números reales?
            
            En Teoría, se ha demostrado que la función generatriz de momentos de una distribución Gamma es:
            \begin{equation*}
                M_X(t) = \left(\dfrac{\lambda}{\lambda - t}\right)^u \qquad t < \lambda
            \end{equation*}

            \item ($1$ punto) A partir de la expresión de la función generatriz de momentos anterior obtener el momento de orden 2 centrado en la media.
            
            Tenemos que nos están pidiendo el momento de orden 2 centrado en la media, es decir, la varianza.
            En Teoría, se ha demostrado que esta es:
            \begin{equation*}
                \text{Var}[X] = \dfrac{u}{\lambda^2}
            \end{equation*}
        \end{enumerate}
    \end{ejercicio}

\end{document}